\documentclass[10pt]{article}
\usepackage{amsmath,amssymb,array}
\usepackage[margin=1in]{geometry}

\title{Summary of the Ternary Quartic Rank-4 Proof Process}
\date{2026-04-13}

\begin{document}
\maketitle

\section*{Scope}
This note summarizes how the theorem
\[
  \texttt{TernaryQuartic.ternaryQuartic\_rankFour\_no\_spurious\_socp}
\]
was found and formalized on the branch
\texttt{autoproof-tq-r4/2026-04-09}. I reviewed the full branch history from
the merge-base \texttt{0c5c820} through the proof-completing commit
\texttt{7d4f423}; the later head \texttt{091ce22} only tweaks scripts. The
branch contains many transport wrappers, coefficient lemmas, and bookkeeping
commits; the list below isolates the ones that changed the proof plan or closed
an open mathematical branch.

\section*{How the Proof Evolved}
The starting point was the abstract SOCP certificate. For a rank-$4$ factor
$u=(u_0,u_1,u_2,u_3)$ and quartic target $p$, the residual is
$r(u,p)=\sigma(u)-p$ with $\sigma(u)=\sum_i u_i^2$. The first decisive idea was
to avoid searching directly for a full quartic identity and instead prove the
vanishing of the residual whenever every SOS square $q^2$ can be decomposed
into an admissible image term plus admissible kernel squares for the linearized
map $A_u(v)=\sum_i u_i v_i$. Commit \texttt{2c3fa6a} introduced exactly this
certificate layer in Lean.

The first Julia pass then changed the proof search from ``look for a counterexample''
to ``classify the only possible quotient obstructions.'' The script
\texttt{julia/tq\_r4\_exploration\_2026\_04\_09.jl} found no numerical
counterexample in random fixed-$u$ searches or in the $15$ monomial-basis
subspaces, and it showed that representative image ranks were always
$13$, $14$, or $15$. More importantly, it displayed the missing monomials and
explicit kernel directions in the mixed-affine models. This is what turned the
project into a branch classification problem rather than a direct dual-certificate
translation.

The next stage solved representative branches one by one. The surjective
constant model $(1,x_0^2,x_0x_1,x_1^2)$ was easy because every quartic monomial
lies in the image. The mixed-affine models of ranks $13$, $14$, and $15$ were
then closed by identifying the missing quotient monomials and matching them
with explicit kernel squares. The crucial proof-serving input was the
six-monomial quadratic normal form from
\texttt{QuadraticNormalForm.lean}, which let Lean read off the required
quartic coefficients of $q^2$ exactly.

At that point the proof was still only representative-level. The branch would
not close without a way to move an arbitrary admissible factor to those solved
models. This is why the transport layer was added next. The files
\texttt{AffineTransform.lean}, \texttt{AffineSocpTransform.lean}, and
\texttt{RepresentativeTransport.lean} formalized two invariances: affine
changes of variables on the polynomial side and orthogonal mixing of the four
factor coordinates on the coefficient-vector side. The point of this layer was
not aesthetic. It separated the actual mathematics into a clean statement:
produce normalization data taking $u$ to a solved representative or solved
plane configuration, and the residual vanishes automatically.

The main middle-period change in strategy was the move from rigid
representatives to basis-invariant plane theorems. The mixed-affine branch
showed that exact representative normalization was often unnecessarily rigid.
Julia experiments on determinant-zero and constant-containing mixed-affine
families indicated that the real invariant was the quadratic plane, or
equivalently the binary quadratic annihilator of that plane. This led to the
completion-of-squares normal forms formalized in
\texttt{MixedAffineNormalization.lean} and then to
\texttt{RepresentativeMixedAffinePlane.lean}, where the representative proofs
were lifted from specific coordinates to orthonormal in-plane relation data.
An analogous move happened for the low-affine branch in
\texttt{RepresentativeLowAffine.lean}. This was the structural step that kept
the final argument manageable.

The last and longest phase was the exact-affine classifier in
\texttt{ExactAffineClassification.lean}. Here the key object is the submodule
of coefficient vectors whose relations have no homogeneous quadratic part:
\[
  \texttt{exactAffineSubmodule}(u)
  = \ker(\texttt{homCoeffMap}(u)).
\]
Once \texttt{relationPolyLin u} is injective, rank-nullity forces the exact
affine relation space to have dimension $1$, $2$, or $3$. Each value leads to
a different geometric model. Dimension $3$ is the span-three branch, where one
can reconstruct $1,x_0,x_1$ and use the abstract span-three theorem. Dimension
$2$ splits into low-affine and mixed-affine branches depending on whether a
constant relation exists. Dimension $1$ is the difficult affine-rank-one
regime, where one normalizes the unique exact affine line and then studies the
remaining quadratic plane by range, discriminant, and support.

The final proof found is therefore not ``decode one SDP dual certificate and
push it into Lean.'' It is a symbolic classification proof guided by Julia but
closed entirely by exact affine and quadratic normal forms in Lean.

\section*{Final Proof Structure}
The root theorem in \texttt{TernaryQuarticProof.lean} has a short statement
because all of the work is hidden in the classifier.

\begin{enumerate}
\item If \texttt{relationPolyLin u} has nontrivial kernel, then the four
coordinates of $u$ satisfy a nontrivial scalar relation
$\sum_i c_i u_i = 0$. For every quadratic factor $q$ in an SOS decomposition of
$p$, the direction $w_i=c_i q$ lies in $\ker A_u$, and the abstract
image-plus-kernel-square certificate gives $r(u,p)=0$. This closes the
dependent case immediately.
\item Assume now that \texttt{relationPolyLin u} is injective. Then the proof
works with \texttt{exactAffineSubmodule(u)}.
\item If the exact-affine dimension is $3$, Lean reconstructs exact relations
for $1,x_0,x_1$. The span-three theorem then shows that every square $q^2$ is
an admissible image term plus a kernel square correction of degree at most $3$,
so the residual vanishes.
\item If the exact-affine dimension is $2$, Lean splits according to whether
the exact affine space contains the constant $1$. Without a constant one lands
in the low-affine plane theorems. With a constant one lands in the mixed-affine
plane theorems. In both cases the image ranks are $13$, $14$, or $15$, and the
missing quotient monomials are matched by explicit kernel families already
formalized at the plane level.
\item If the exact-affine dimension is $1$, Lean normalizes the unique affine
line and analyzes the complementary quadratic plane. This produces the long
affine-rank-one branch decomposition: repeated-line, common-tail, diagonal,
shared-tail, determinant-zero, and range-two subcases. These were resolved
through the combined infrastructure in
\texttt{RepresentativeAffineRankOne.lean} and
\texttt{ExactAffineClassification.lean}.
\item The last unresolved subcase was the constant-containing exact-affine
dimension-$1$ branch. Commit \texttt{7d4f423} closed it by proving that once
the complementary quadratic tail matrix is invertible, one can reconstruct
exact relations for $x_0^2$, $x_0x_1$, and $x_1^2$ from arbitrary independent
quadratic relations. That reduces the branch to the surjective theorem in
\texttt{RepresentativeSurjective.lean}.
\end{enumerate}

This is why the final root proof is short: first split on whether
\texttt{relationPolyLin u} has kernel, and in the injective case call the exact
affine classifier.

\section*{Proof-Defining Commits}
The following commits are the ones that materially advanced the proof. I omit
cleanup commits, documentation-only tweaks, and the many local wrapper commits
whose purpose was to make an already chosen strategy formalizable.

\small
\begin{itemize}
\item \texttt{2c3fa6a}, \texttt{6201318}, \texttt{402ec5f}: created the
abstract certificate layer, recorded the first dual infeasibility evidence, and
closed the dependent case while revealing the first rank-based split.
\item \texttt{bc0a23e}: formalized orthogonal factor mixing and opened the
mixed-affine representative program.
\item \texttt{b45442a}, \texttt{47b10fc}, and the representative cluster ending
at \texttt{2c74558}: supplied the surjective constant branch, the six-monomial
quadratic normal form, and the rank-$13/14/15$ mixed-affine representatives.
\item \texttt{a702d77} and \texttt{86368c4}: solved the span-three branch first
at a representative and then in basis-free form. This removed one whole
normalization problem from the final proof.
\item \texttt{db938e8}, \texttt{1840fb1}, \texttt{355b7ff}: built the affine
transport and representative pullback layer.
\item \texttt{b5b6e37}: turned the mixed-affine Julia evidence into exact
completion-of-squares normal forms.
\item \texttt{c69c9b2} and \texttt{9c92cff}: replaced rigid representatives by
basis-invariant plane theorems for the mixed-affine and low-affine branches.
\item \texttt{0965a6f}: introduced the exact-affine classifier, which became
the backbone of the final theorem.
\item \texttt{77eed41}, \texttt{7d08f55}, and \texttt{50364e9}: closed the
dimension-$2$ no-constant branch and reduced the dimension-$1$ branch to a
short list of affine-rank-one and range-two endpoints.
\item \texttt{7d4f423}: closed the last constant-containing dimension-$1$ gap,
added the final surjective reconstruction lemmas, and assembled the root
theorem.
\end{itemize}
\normalsize

\section*{Role of the Julia Scripts}
The Julia code was useful, but useful in a specific way. It did \emph{not}
produce a numerical certificate that was later imported or rationalized inside
Lean. The final proof is exact and symbolic. What the scripts provided was:

\begin{enumerate}
\item early negative evidence against counterexamples;
\item image-rank and missing-monomial data that exposed the correct quotient
obstructions;
\item explicit candidate kernel families in the representative branches;
\item branch pruning for mixed-affine and exact-affine normal forms.
\end{enumerate}

The main script, \texttt{tq\_r4\_exploration\_2026\_04\_09.jl}, changed the
strategy outright: it showed that all monomial-basis cases looked infeasible,
that the representative image ranks were only $13$, $14$, or $15$, and that
the mixed-affine missing classes were exactly the monomials corrected later by
Lean kernel families. The 2026-04-10 scripts were the next most useful. They
tested determinant-zero and constant-perturbed families using exact coefficient
linear algebra and showed that only a small number of annihilator normal forms
were genuinely nonsurjective. The 2026-04-11 and 2026-04-12 scripts for the
affine-dimension-$1$ branch were still valuable, but mainly as branch-pruning
tools rather than sources of directly formalizable certificates.

So the correct assessment is that Julia supplied strong strategic guidance,
especially in identifying the right invariants, normal forms, and kernel
families, but not the final proof artifact. The final theorem is an exact Lean
classification whose shape was suggested, rather than certified, by the
experiments.

\end{document}
